\section{SVD、低秩与数值秩}
\label{sec:08-Numerical-Analysis/05-Numerical-Linear-Algebra/06}

一个 \(4096\times4096\) 的矩阵，可能是某一层的权重，也可能是一批样本的嵌入。你调 \texttt{rank}，它回答 4096，满秩。你转手做一次 PCA，前 60 个方向解释了 \(95\%\) 的方差。同事换了个容差重跑，\texttt{rank} 这次回答 212。

三个数字都没错，错的是以为``秩''在浮点世界里仍是一个离散的、只取决于矩阵本身的整数。上一节的特征值至少还有 Weyl 定理兜底，秩连这个都没有——它是一个先落在整数上、再被舍入误差抹平的量。

\subsection{奇异值是矩阵的一列放大倍数}

奇异值分解把任意实矩阵写成

\[
A=U\Sigma V^{\trans},\qquad
\Sigma=\diag(\sigma_1,\ldots,\sigma_p),\quad
\sigma_1\ge\cdots\ge\sigma_p\ge0
\]

（复情形把转置换成共轭转置）。\(V\) 的第 \(i\) 列 \(v_i\) 被 \(A\) 放大 \(\sigma_i\) 倍并转到 \(u_i\) 上，即 \(Av_i=\sigma_iu_i\)。\hyperref[sec:03-Linear-Algebra/07-SVD-and-Data-Applications/01]{``任意矩阵的奇异值分解''}已经把这套几何讲透了，这里只用它的一个推论：矩阵的全部信息，都编码在 \(\sigma_i\) 这列数从大到小怎么衰减。秩是这列数里非零的个数，条件数是首尾之比，零空间是尾部方向张成的空间——三件事共用同一份数据。

计算上有一条必须守住的纪律：不要先形成 \(A^{\trans}A\) 再求特征值。它把条件数平方，\(\kappa_2(A^{\trans}A)=\kappa_2(A)^2\)，本来还能分辨的小奇异值平方之后直接沉到舍入尺度以下。标准算法先用正交变换把 \(A\) 双对角化，再在双对角矩阵上迭代，全程不显式构造 Gram 矩阵。

\subsection{奇异值谱上找不到一条零线}

后向稳定的 SVD 算法算出的不是 \(A\) 的奇异值，而是某个 \(A+E\) 的精确奇异值，其中 \(\norm{E}_2\le c(n)\,\varepsilon\,\norm{A}_2\)。配上 Weyl 扰动界 \(\abs{\hat\sigma_i-\sigma_i}\le\norm{E}_2\)，得到

\[
\abs{\hat\sigma_i-\sigma_i}\;\lesssim\;c(n)\,\varepsilon\,\sigma_1 .
\]

右端与 \(i\) 无关。这就是全部麻烦的来源：算法对每个奇异值的绝对分辨率都是同一个数，约等于 \(\sigma_1\) 乘机器精度再乘一个尺寸因子。理论上非零但落在这条线以下的奇异值，与真正的零在数值上无法区分——不是算得不够好，是这个精度下没有信息可用。

于是实践里换一个定义：数值秩是阈值以上的奇异值个数，

\[
\rank_\tau(A)=\#\set{i:\sigma_i>\tau},\qquad
\tau=\sigma_{\max}\cdot n\cdot\varepsilon .
\]

\begin{center}
\begin{tikzpicture}[x=0.36cm,y=0.21cm,font=\small,>=stealth]
  \draw[black!55] (0.35,-17) -- (25.2,-17);
  \draw[black!55] (0.35,-17) -- (0.35,1.2);
  \foreach \y/\lab in {0/{\(\sigma_1\)},-4/{\(10^{-4}\sigma_1\)},-8/{\(10^{-8}\sigma_1\)},
                       -12/{\(10^{-12}\sigma_1\)},-16/{\(10^{-16}\sigma_1\)}}{
    \draw[black!45] (0.05,\y) -- (0.35,\y);
    \node[left,font=\scriptsize,black!65] at (0.0,\y) {\lab};}
  \foreach \i/\v in {1/0, 2/-0.25, 3/-0.45, 4/-0.7, 5/-0.95, 6/-1.25, 7/-1.6, 8/-2.0,
                     9/-2.5, 10/-3.1, 11/-3.9, 12/-4.9, 13/-6.2, 14/-7.8, 15/-9.6,
                     16/-11.5, 17/-13.2}{
    \fill[blue!58] (\i-0.32,-17) rectangle (\i+0.32,\v);}
  \foreach \i/\v in {18/-14.6, 19/-15.4, 20/-15.8, 21/-15.9, 22/-16.0, 23/-15.95,
                     24/-16.05, 25/-15.85}{
    \fill[black!32] (\i-0.32,-17) rectangle (\i+0.32,\v);}
  \draw[dashed,red!70!black,thick] (0.35,-14.3) -- (25.6,-14.3);
  \node[right,font=\scriptsize,red!70!black] at (18.6,-12.6) {\(\tau=\sigma_{\max}n\varepsilon\)};
  \draw[dotted,black!60] (0.35,-16.0) -- (25.6,-16.0);
  \node[right,font=\scriptsize,black!60] at (25.6,-16.0) {舍入地板};
  \draw[dashed,black!55] (17.5,-17) -- (17.5,-1.5);
  \node[above,font=\scriptsize,black!70] at (14.0,-1.6) {数值秩 \(=17\)};
  \node[below,font=\footnotesize,black!70] at (13,-18.6) {奇异值序号 \(i\)};
\end{tikzpicture}

\emph{纵轴是对数刻度。前十七个奇异值跨越十三个数量级仍在稳步衰减，之后的一段忽然平掉——那不是矩阵的结构，是算法的分辨率下限；阈值线画在哪里，报出来的秩就是多少。}
\end{center}

图里那段平台解释了开头的三个数字。\texttt{numpy.linalg.matrix\_rank} 的默认阈值正是 \(\sigma_{\max}\cdot\max(m,n)\cdot\varepsilon\)，对应图中那条红线，报 17；把阈值调低到地板附近，平台上的每一根柱子都会被算进去，报 25；而 PCA 关心的是能量占比，它的截断点由 \(\sum_{i\le k}\sigma_i^2\big/\sum_i\sigma_i^2\) 决定，落在图的左半边，报的数还要小得多。三个答案对应三条不同的横线，问``这个矩阵的秩是几''而不说横线画在哪，问题本身就没定义完。

阈值还得随场景往上抬。舍入误差只是地板的一种成因；数据本身若带测量噪声，噪声引入的扰动 \(\norm{E}_2\) 可能远大于 \(\varepsilon\sigma_1\)，此时该用噪声水平而不是机器精度当阈值。另一件容易忽略的事是尺度：把某几列换个单位，全部奇异值会整体缩放甚至改变相对大小，而``秩为 7''这句话跟着一起变。没有说明尺度和容差的秩，不是一个完整的数值结论。

\subsection{截断点该选在哪里}

保留前 \(k\) 项得到 \(A_k=U_k\Sigma_kV_k^{\trans}\)。Eckart--Young 定理说它在所有秩不超过 \(k\) 的矩阵里最优：

\[
\norm{A-A_k}_2=\sigma_{k+1},
\qquad
\norm{A-A_k}_F=\Bigl(\sum_{i>k}\sigma_i^2\Bigr)^{1/2}.
\]

骨架是这样的：任何秩 \(k\) 矩阵 \(B\) 的零空间维数至少 \(n-k\)，它必与 \(v_1,\ldots,v_{k+1}\) 张成的 \(k+1\) 维空间相交于某个非零单位向量 \(w\)；在这个 \(w\) 上 \(Bw=0\)，于是 \(\norm{A-B}_2\ge\norm{Aw}_2\ge\sigma_{k+1}\)。截断 SVD 恰好取到这个下界。条件对账：结论对 2-范数与 Frobenius 范数成立，对一般范数不自动成立；\(\sigma_k=\sigma_{k+1}\) 时最优解不唯一，截断只是其中一个；这里说的最优是对给定的 \(A\) 而言，若 \(A\) 本身含噪，被逼近得最好的是含噪的那份数据。

于是截断点怎么选就有了两种情形。谱上有明显间隙时，选在间隙处，此时 \(k\) 的小幅变动几乎不影响结果，结论稳健。谱平滑衰减时（真实数据的常态），间隙不存在，\(k\) 只能由目标定：要控制重构误差就按 \(\sigma_{k+1}\) 反解，要控制解释方差就按能量比反解，要控制下游任务指标就只能扫一遍 \(k\) 拿验证集说话。把``秩''当成需要发现的客观事实去找，找不到；把它当成一个需要声明的建模选择，就清楚了。

这也解释了低秩适配里那个反复出现的问题：\(r\) 该取多少。权重矩阵的数值秩通常接近满秩，它并不直接给出答案，因为被拟合的是更新量而不是权重本身；有用的是谱衰减的形状——衰减快意味着少数方向已经装下了大部分变化，衰减慢意味着截断一定会丢东西，代价得提前算进预算。\hyperref[sec:03-Linear-Algebra/07-SVD-and-Data-Applications/03]{``低秩近似：用少数方向保留主要结构''}讲的是这件事的代数一面，这里补的是它的数值一面：在浮点下，那条曲线的尾部有一段是看不见的。

\subsection{小奇异值方向上的解不值得相信}

伪逆 \(A^+=V\Sigma^+U^{\trans}\) 把非零奇异值取倒数。落在地板附近的 \(\sigma_i\) 一旦被取倒数，就成了 \(10^{16}\) 量级的放大器，而它乘的恰好是那个方向上的噪声。最小二乘解在这个方向上的分量因此完全由噪声决定，范数还可能大得离谱。

补救办法都是同一个形状：衰减而不是求逆。截断 SVD 把阈值以下的方向直接丢掉，Tikhonov 正则化把 \(1/\sigma_i\) 换成 \(\sigma_i/(\sigma_i^2+\mu)\)，两者都在小 \(\sigma_i\) 处平滑地趋零。代价是引入偏差，而 \(\mu\) 或截断点就是偏差与方差之间的那个旋钮——\hyperref[sec:03-Linear-Algebra/07-SVD-and-Data-Applications/02]{``伪逆与病态问题''}和\hyperref[sec:08-Numerical-Analysis/01-Floating-Point-and-Error/03]{``条件数衡量问题本身的敏感性''}从两个方向讲的是同一件事。

验证 SVD 结果时也要挑对该看的量。奇异值本身稳定，可以逐个比较；奇异向量在奇异值接近时可以在对应子空间里大幅旋转，符号更是本来就不唯一，所以该比较的是子空间夹角而不是向量分量。可算的健全性检查有三个：重构残差 \(\norm{A-U\Sigma V^{\trans}}_2\) 是否在 \(\varepsilon\norm{A}_2\) 量级，\(U\) 与 \(V\) 的正交性偏差是否同样量级，奇异值是否非负且有序。

矩阵大到装不下时，只需要前 \(k\) 个方向的话，完整 SVD 是浪费。Lanczos 双对角化只用矩阵--向量乘，随机化 SVD 先用随机投影捕获近似列空间再在小矩阵上做精确分解——后者的精度取决于过采样量、幂迭代次数和谱衰减速度，输出形状虽与 SVD 相同，报告时必须附上残差和随机参数，不能当成精确分解使用。

不过上面所有做法都还共享一个前提：矩阵能整个拿到手，哪怕只是以``能乘一次''的方式。下一节把这个前提也拆掉，看看当内存一次只装得下一行时，线性系统还能怎么解。
